\chapter{Implementierung des Backtesters und der Strategien}
\label{cha:ImplementierungdesBacktestersundderStrategien}

In diesem Kapitel wird beschrieben wie die einzelnen Komponenten des Backtester und die getesteten Strategien implementiert sind.
\iffalse
\begin{CCode}[caption={Simple C example}, label={lst:c-example}]
#include <stdio.h>

int main(void)
{
    printf("Hello, world!\n");
    return 0;
}
\end{CCode}

\begin{HtmlCode}[caption={Basic HTML page}, label={lst:html-example}]
<!DOCTYPE html>
<html lang="en">
<head>
    <meta charset="UTF-8">
    <title>Example</title>
</head>
<body>
    <h1>Hello World</h1>
    <p>This is an HTML example.</p>
</body>
</html>
\end{HtmlCode}

\begin{JsCode}[caption={JavaScript example}, label={lst:js-example}]
function greet(name) {
    console.log(`Hello, ${name}!`);
}

greet("Alice");
\end{JsCode}
\fi

\section{Die Benutzerschnittstellen}
\label{sec:DieBenutzerschnittstellen}

\subsection{Die API Schnittstelle}
\label{subsec:DieAPISchnittstelle}

\subsection{Die Logikeingangsparameter}
\label{DieLogikeingangsparameter}

\section{Implementierung der Backtestlogik}
\label{sec:ImplementierungenderBacktestlogik}

\section{Implementierungen der Strategien}
\label{sec:ImplementierungenderStrategien}

Es werden drei verschiedene Handelsstrategien implementiert. Jede Strategie hat eine unterschiedliche Logik, um Kauf- und Verkaufsentscheidungen zu treffen.
Diese Unterschiede in der Logik führen zu unterschiedlichen Laufzeitkomplexitäten.
Die Ideen hinter den implementierten Strategien sind in Kapitel \ref{sec:Bewährte AnsätzefürHandelsstrategien} beschrieben.

Jede der folgenden Untersektionen beschreibt die Implementierung einer der drei Strategien in JavaScript und C.

\subsection{Gleitende Durchschnittsüberschneidungs Strategie}
\label{subsec:GleitendeDurchschnittsüberschneidungsStrategie}

\begin{CCode}[caption={Durchschnittsüberschneidung in C}, label={lst:DurchschnittsüberschneidunginC}]
    for (int i = prices_length - 1 - shortPeriod; i < prices_length - 1; i++) {
        shortPrev += prices[i];
    }
    shortPrev /= shortPeriod;

    for (int i = prices_length - 1 - longPeriod; i < prices_length - 1; i++) {
        longPrev += prices[i];
    }
    longPrev /= longPeriod;

    double shortCurr = 0.0;
    double longCurr = 0.0;

    for (int i = prices_length - shortPeriod; i < prices_length; i++) {
        shortCurr += prices[i];
    }
    shortCurr /= shortPeriod;

    for (int i = prices_length - longPeriod; i < prices_length; i++) {
        longCurr += prices[i];
    }
    longCurr /= longPeriod;

    const double diffPrev = shortPrev - longPrev;
    const double diffCurr = shortCurr - longCurr;

    double absDiff = diffCurr < 0.0 ? -diffCurr : diffCurr;

    const double currentPrice = prices[prices_length - 1];
    const double threshold = currentPrice * 0.002;

    if (diffPrev < 0.0 && diffCurr > 0.0 && absDiff >= threshold) {
        return 1;
    }
    if (diffPrev > 0.0 && diffCurr < 0.0 && absDiff >= threshold) {
        return -1;
    }
\end{CCode}

\subsection{Mittelwertrückkehr Strategie}
\label{subsec:MittelwertrückkehrStrategie}

\begin{CCode}[caption={Mittelwertrückkehr in C}, label={lst:MittelwertrückkehrinC}]
    double upperBand = mean + 3.25 * stdDev;
    double lowerBand = mean - 3.25 * stdDev;

    double lastPrice = prices[length - 1];

    if (lastPrice < lowerBand) {
        return 1;
    }
    if (lastPrice > upperBand) {
        return -1;
    }
\end{CCode}

\subsection{Ausbruchsstrategie}
\label{subsec:Ausbruchsstrategie}

\begin{CCode}[caption={Ausbruchsstrategie in C}, label={lst:AusbruchsstrategieinC}]
    double lastPrice = prices[length - 1];
    double epsilon = 0.01;

    double longTrigger = max * (1.0 + epsilon);
    double shortTrigger = min * (1.0 - epsilon);

    if (lastPrice > longTrigger) {
        return 1;
    }
    if (lastPrice < shortTrigger) {
        return -1;
    }
\end{CCode}

\section{Die Datenvisualisierungen}
\label{sec:DieDatenvisualisierungen}

\subsection{Visualisierung des Kurswertes}
\label{subsec:VisualisierungdesKurswertes}

\subsection{Visualisierung des Portfoliowertes}
\label{subsec:VisualisierungdesPortfoliowertes}